\begin{abstract}
\noindent 
The pricing kernel — the stochastic discount factor linking risk and return across all states of the world — simultaneously determines fair derivative prices, risk premia magnitudes, and optimal hedge ratios, making it the central object of modern asset pricing. For Bitcoin, the world's largest cryptocurrency with a derivatives market exceeding \$40 billion in open interest, its shape and dynamics remain largely uncharacterized. Existing Bitcoin option pricing studies operate under either the physical measure or the risk-neutral measure in isolation, and no study has yet imposed a unified Stochastic Volatility with Correlated Jumps (SVCJ) framework simultaneously across both measures, leaving recovered pricing kernels internally inconsistent and the Bitcoin pricing kernel puzzle unresolved. This thesis estimates both measures within a common SVCJ specification — the physical measure via rolling MCMC on spot returns and the risk-neutral measure via daily Levenberg–Marquardt calibration on Deribit options — recovering the pricing kernel as an internally consistent density ratio and deploying it as a trading signal to test whether Bitcoin pricing kernel anomalies are economically exploitable. The study covers January 2020 to December 2024, using daily Bitcoin spot returns and a comprehensive panel of European-style inverse options across multiple strikes and maturities from Deribit. The analysis proceeds through four stages: rolling MCMC estimation under $\mathbb{P}$
, daily calibration under $\mathbb{Q}$
, nonparametric pricing kernel estimation with monotonicity tests, and backtesting of a Variance Risk Premium-filtered delta-hedged short straddle strategy.
\vspace{0in}\\
\noindent\textbf{Keywords:} Bitcoin, pricing kernels, stochastic discount factor, stochastic volatility with correlated jumps, cryptocurrency derivatives\\
\vspace{0in}\\
\noindent\textbf{JEL Codes:} JEL1, JEL2, JEL3\\

\bigskip
\end{abstract}
